\subsection*{頭字語}
\addcontentsline{toc}{subsection}{頭字語}
\par\smallskip\noindent\refstepcounter{table}\label{tab:ch15-02}表 \thetable: 頭字語の整理\par\smallskip
表 \thetable は、頭字語の整理を比較・確認しやすい形で整理したものである。\par\smallskip
\begin{longtable}{L{0.20\linewidth} L{0.34\linewidth} L{0.40\linewidth}}
\toprule
略語 & 英語 & 日本語での意味 \\
\midrule
IRR & Internal Rate of Return & 内部収益率。投資案の収益性を示す利率である。 \\
MARR & Minimum Acceptable Rate of Return & 許容最低収益率。組織が投資として受け入れる最低限の収益率である。 \\
SDLC & Software Development Life Cycle & ソフトウェア開発ライフサイクル。開発活動の流れである。 \\
SPLC & Software Product Life Cycle & ソフトウェア製品ライフサイクル。企画、開発、運用、保守、廃止までを含む。 \\
ROI & Return on Investment & 投資収益率。投資に対する利益の割合である。 \\
SEI & Software Engineering Institute & ソフトウェア工学研究所。 \\
TCO & Total Cost of Ownership & 総所有費用。取得から運用・保守までの総費用である。 \\
\bottomrule
\end{longtable}
